\documentclass[]{report}


%%%%%%%%%%%%%%%%%%%%%%%%%%%%%%%%%%%%%%%%%%%%%%%%%%%%%%%%%%%%%%%%%%%%%%%%%%%%%%%%%%%%%%%
% Pakete und Einstellungen
%%%%%%%%%%%%%%%%%%%%%%%%%%%%%%%%%%%%%%%%%%%%%%%%%%%%%%%%%%%%%%%%%%%%%%%%%%%%%%%%%%%%%%%

% The cmap package is intended to make the PDF files generated by 
% pdflatex "searchable and copyable" in acrobat reader and other
% compliant PDF viewers.
\usepackage{cmap}


% Verwendung des UTF-8 Zeichensatzes mit Umlauten
\usepackage[utf8]{inputenc}

% Um Fonts in westeuropäischer Codierung zu verlangen, verbessertes Suchen und Kopieren
% aus PDF-Dokumenten
\usepackage[T1]{fontenc} 

%Schriftart Latin Modern statt Computer Modern (weniger pixelig, erscheint flüssiger)
% als Standard-Serifen-Schriftart
\usepackage{lmodern}

% Helvetica (resp. Open-Source-Variante) als Standard-Sans-Serif-Schriftart
\usepackage[scaled]{helvet}


% Verwende Standardmässig für gesamtes Dokument die Standard-Sans-Serif-Schriftart
% Auskommentieren für Serifen-Schriftart
\renewcommand{\familydefault}{\sfdefault}

% Indexe in Inhaltsverzeichnis anzeigen
\usepackage[]{tocbibind}

% English
\usepackage[english]{babel}

% Um Sachen einzufärben
\usepackage{color}
\usepackage{xcolor}


% Einbinden von Grafiken und Bildern
\usepackage{graphicx}

\graphicspath{ {./assets/} }
\usepackage{subfig}

% Für Tabellennotizen
\usepackage{tablefootnote}
\usepackage{pdflscape}
\usepackage{tikz}



% Einbinden von Code
\usepackage{listings}
% Styling des Codebereiches
\definecolor{codegreen}{rgb}{0,0.6,0}
\definecolor{codegray}{rgb}{0.5,0.5,0.5}
\definecolor{codepurple}{rgb}{0.58,0,0.82}
\definecolor{darkblue}{rgb}{0.1, 0.1, 0.4}
\lstdefinestyle{mystyle}{   
	commentstyle=\color{codegreen},
	keywordstyle=\color{magenta},
	numberstyle=\tiny\color{codegray},
	stringstyle=\color{codepurple},
	basicstyle=\ttfamily\footnotesize,
	breakatwhitespace=false,         
	breaklines=true,                 
	captionpos=b,                    
	keepspaces=true,                 
	numbers=left,                    
	numbersep=5pt,                  
	showspaces=false,                
	showstringspaces=false,
	showtabs=false,                  
	tabsize=2
}
\lstset{style=mystyle}

% for nicer formulas
\usepackage{amsmath}

% Das Setspace Paket ermöglicht es auch recht einfach Weise den Zeilenabstand zu ändern
\usepackage[onehalfspacing]{setspace}

% Paket zur Gestaltung und Erstellung von Literaturverzeichnissen und Zitaten
\usepackage{csquotes}
\usepackage[bibencoding=utf8, backend=biber, style=apa, dashed=false, maxbibnames=99, doi=true, isbn=false]{biblatex}
\addbibresource{pointclouds.bib}
\DefineBibliographyStrings{english}{andothers = {et al.}}

% Damit die URL in der Literatur an mehr Orten Umgebrochen wird
\setcounter{biburllcpenalty}{7000}
\setcounter{biburlucpenalty}{8000}

% Überschreibung der Standard-Titel-Schriftart bei Thesis und Incollection,
% damit auch diese im Literaturverzeichnis kursiv gargestellt werden.
\DeclareFieldFormat[thesis]{title}{\textit{{#1}}}
\DeclareFieldFormat[incollection]{title}{\textit{{#1}}}

% Setze Sortierung der Namen auf Nachname, Vorname
\DeclareNameAlias{sortname}{family-given}

%\DeclareDelimFormat{multinamedelim}{\addspace\slash\space}
%\DeclareDelimAlias{finalnamedelim}{multinamedelim}

% damit die Einträge in der Bibliographie nicht aneinander kleben
\setlength\bibitemsep{1.5\itemsep}

% Paket, um die Seite zu gestalten
\usepackage[a4paper,left=25mm,right=36mm,top=21mm,bottom=27mm,footskip=10mm]{geometry}

% Paket zur einfacheren Manipulation von Kopf- und Fußzeile
\usepackage{fancyhdr}


\fancypagestyle{plain}{
	\fancyhf{} % clear all header and footer fields
	%\fancyhead[L]{\includegraphics[height=0.8cm]{assets/FHNW_HABG_10mm.jpg}}
	\fancyfoot[R]{\textsf{\thepage}} % except the center
	\fancyfoot[L]{\textsf{FHNW, HABG, Institute Geomatics}}
	\renewcommand{\headrulewidth}{0pt}
	\renewcommand{\footrulewidth}{1pt}
}

\fancypagestyle{titlepage}{
	\fancyhf{} % clear all header and footer fields
	\fancyhead[L]{\hspace*{-0.7cm}
	\setlength{\headheight}{33pt}
	\includegraphics[height=1cm]{assets/FHNW_en.PNG}}
	\renewcommand{\headrulewidth}{0pt}
	\renewcommand{\footrulewidth}{0pt}
}

\usepackage{hyperref}


\usepackage{caption}
\captionsetup[figure]{name={Img.}}

% Damit Sonderzeichen über Befehle wie \textdegree funktionieren
\usepackage{textcomp}

%für bilder in zwei kolonnen
\usepackage{wrapfig}

%um PDFs als Anhang zu integrieren
\usepackage{pdfpages}

\usepackage{multicol}

% für enumerate roman stuff
\usepackage{enumitem}

% um grosse teile auszukommentieren
\usepackage{comment}

\usepackage{float}
\usepackage{booktabs}
\usepackage{changepage}
\usepackage{pgfgantt}


%%%%%%%%%%%%%%%%%%%%%%%%%%%%%%%%%%%%%%%%%%%%%%%%%%%%%%%%%%%%%%%%%%%%%%%%%%%%%%%%%%%%%%%
% Dokument
%%%%%%%%%%%%%%%%%%%%%%%%%%%%%%%%%%%%%%%%%%%%%%%%%%%%%%%%%%%%%%%%%%%%%%%%%%%%%%%%%%%%%%%

\begin{document}
	\pagestyle{plain}
	\setlength\parindent{0pt}
	
	
	%%%%%%%%%%%%%%%%%%%%%%%%%%%%%%%%%%%%%%%%%%%%%%%%%%%%%%%%%%%%%%%%%%%%%%%%%%%%%%%%%%%%%%%
	% Titelseite
	%%%%%%%%%%%%%%%%%%%%%%%%%%%%%%%%%%%%%%%%%%%%%%%%%%%%%%%%%%%%%%%%%%%%%%%%%%%%%%%%%%%%%%%
	
	\begin{titlepage}
		\thispagestyle{titlepage}
		\pagenumbering{Roman} 
		\setcounter{page}{1}
		\vspace*{1cm}
		
		\begin{flushleft}
			
			{\Large HABG VP2}\\
			\vspace*{1mm}
			{\Large MSE: Master of Science in Engineering}
			\vspace*{3cm}
			
			
			{\huge \textbf{{VP2: Integration and Visualization of 3D Data
for the Ballenberg Open-Air Museum}}}\\
			
			
			\begin{figure}[h]
            \centering
            \includegraphics[height = 8cm]{assets/placeholder_title.png}
            \caption{Potree Ballenberg IGEO FHNW}
            \end{figure}
			
			\vfill
			
			{\Large
				\begin{tabbing}
					\hspace*{3.5cm}\= \kill
					Autor: \> Joël Bachmann \\
					Supervisors: \> Prof. Dr. Wissam Wahbeh \\
                                \> Prof. Dr. David Grimm \\
                    Expert:      \> 
					\\ \\
				\end{tabbing}
				
				Baden, 11.03.2026
			}
		

			
			
		\end{flushleft}
		
		
	\end{titlepage}
	
	
	%%%%%%%%%%%%%%%%%%%%%%%%%%%%%%%%%%%%%%%%%%%%%%%%%%%%%%%%%%%%%%%%%%%%%%%%%%%%%%%%%%%%%%%
	% Zusammenfassung
	%%%%%%%%%%%%%%%%%%%%%%%%%%%%%%%%%%%%%%%%%%%%%%%%%%%%%%%%%%%%%%%%%%%%%%%%%%%%%%%%%%%%%%%
	
	
	\begin{titlepage}
	    
		\thispagestyle{plain}
		\pagenumbering{Roman} 
		\setcounter{page}{2}
		
% No TOC entry here — will be done below
\begin{abstract}
Lorem Ipsum
\end{abstract}




	
		
		
	\end{titlepage}

\setcounter{page}{3}

\addtocontents{toc}{\protect\setcounter{tocdepth}{-1}}
\tableofcontents
\addtocontents{toc}{\protect\setcounter{tocdepth}{3}}




	
	%%%%%%%%%%%%%%%%%%%%%%%%%%%%%%%%%%%%%%%%%%%%%%%%%%%%%%%%%%%%%%%%%%%%%%%%%%%%%%%%%%%%%%%
	% Arbeit
	%%%%%%%%%%%%%%%%%%%%%%%%%%%%%%%%%%%%%%%%%%%%%%%%%%%%%%%%%%%%%%%%%%%%%%%%%%%%%%%%%%%%%%%
	\newpage
	\pagenumbering{arabic} 


\chapter{State of the Art}
\section{Introduction}
\label{sec:introduction}

The digital documentation of built cultural heritage has matured rapidly over the past two
decades. Terrestrial laser scanning, UAV photogrammetry, and close-range photogrammetry can
now produce geometrically accurate, high-resolution three-dimensional records of historic
buildings at a fraction of the cost and time that equivalent measured surveys once required.
Many institutions have seized this opportunity and built substantial 3D archives. The
challenge that now confronts them is no longer acquisition but management: how to organize,
describe, and deliver datasets that were captured at different times, with different
instruments, in different formats, and by different teams --- and how to make this
heterogeneous archive usable not only for the surveyors who created it, but for conservators,
researchers, curators, and the visiting public.

The Ballenberg Open-Air Museum in the Swiss canton of Bern is a representative and
particularly demanding example of this challenge. The museum preserves more than 100
historic rural buildings relocated from across Switzerland, representing vernacular
construction traditions from every region of the country and spanning several centuries of
building practice. Documentation campaigns have been conducted over multiple years, using
terrestrial laser scanning and photogrammetric methods, producing a growing archive of
heterogeneous 3D data stored at \texttt{S:\textbackslash HABG\textbackslash
A1334\_IGEO\_Projekte\textbackslash A1334\_Projekte\_divers\textbackslash 2025\textbackslash
Ballenberg}. No unified pipeline currently exists to integrate these datasets, attach
structured metadata, or deliver them to web-based viewers accessible to the museum's
diverse user groups. The project described in this report addresses precisely these gaps.

Three potential user groups are hypothesized for a future Ballenberg platform, pending
validation through user research. Museum staff --- conservators and curators with domain
knowledge but limited point cloud experience --- would benefit from measurement tools and
rapid access to condition data for maintenance planning. A second group, museum visitors
including those with restricted physical access to individual buildings, would require
simple, focused views that highlight a single architectural or historical aspect rather
than an open-ended navigable point cloud. A third hypothesized group consists of students
participating in guided educational workshops, who would engage with the collection through
narrative-driven explorations tied to on-site visits. These three groups inform the
platform and interaction requirements discussed in Section~\ref{sec:visualization},
particularly the emphasis on mobile optimization and minimal interaction complexity for
non-specialist audiences.

Before any such pipeline can be designed, the available technical options must be understood.
This State of the Art report reviews the current literature and practice in four
interconnected themes: the formats and standards used to store and exchange point cloud
data; the strategies for structuring and indexing heterogeneous 3D archives; the use of
semantic metadata to attach meaning to 3D geometry; and the open-source web platforms
available for delivering large point clouds and 3D models to browser-based audiences. These four themes were chosen because they
represent the principal decision points in any pipeline from raw survey data to accessible
web visualization, and because the state of practice in each directly constrains what is
achievable in the others.

The review is guided by five research questions:

\begin{enumerate}
  \item What data structuring and metadata strategy enables scalable, interoperable, and
        future-proof management of heterogeneous cultural heritage point cloud datasets?
  \item Which web-based platform or combination of platforms provides the best trade-off between
        usability, performance, and functionality for the different user personas of an
        open-air museum?
  \item How can heterogeneous 3D survey products --- terrestrial laser scanning,
        photogrammetry, and mesh reconstructions --- be standardized into a unified
        pipeline for web deployment?
  \item How can semantic metadata be effectively linked to 3D point cloud visualization
        to support documentation, conservation, and research workflows?
  \item Which available Ballenberg datasets are best suited to serve as mobile-optimized
        visualization prototypes, and how can they be used to illustrate the architectural
        and historical dimensions of the collection to the anticipated user groups?
\end{enumerate}


\section{Point Cloud Formats and Standards}
\label{sec:formats}

Point cloud data can be stored in a wide range of formats, each developed with different goals
in mind --- some focus on interoperability, others on performance, compression, or long-term
storage. Choosing the right format matters because it affects how data can be processed, shared,
and accessed in the future. This section gives an overview of the most relevant formats for a
heritage documentation project like Ballenberg.

\subsection{LAS and LAZ}

LAS is a binary format developed and maintained by the American Society for Photogrammetry
and Remote Sensing (ASPRS). It was designed to store LiDAR point cloud data, combining GPS,
IMU, and laser pulse range data into X, Y, Z point records \cite{asprs_las_spec}. Each point can
also store attributes like intensity, return number, and surface classification. LAS~1.4, the
current version, was approved in 2011 with the most recent revision in 2019, and uses a
structured binary layout with a public header block and variable length records \cite{asprs_las_spec}.
LAZ is simply a compressed version of LAS. It achieves a typical compression ratio of 7:1 to
10:1, and the open-source LASzip library is built into many point cloud tools, making it widely
compatible while cutting file sizes \cite{cadinterop_las}. LAS and LAZ are the standard in
geospatial and airborne survey work. However, they have a notable limitation for building
documentation: they cannot store the panoramic images and scanner position metadata that are
typically captured alongside terrestrial scans.

\subsection{E57}

E57 is an open standard developed by the ASTM E57 Committee, designed specifically to
exchange data between different scanning hardware and software. The format stores 3D point
cloud data, associated 2D imagery, and core metadata, with its requirements shaped by consensus
among laser scanner manufacturers, software vendors, and academic researchers \cite{huber2011e57}.
Compared to LAS, E57 is better suited to terrestrial scanning because it packs everything into
one self-contained file, including scanner station positions, calibration data, and panoramic
images. For long-term archiving, E57 is a strong choice because the specification is public and
managed by an independent international body, making it less vulnerable to any single vendor
discontinuing support \cite{cadinterop_e57}. It is worth noting that no manufacturer uses E57 as
their native recording format --- each vendor records in their own proprietary format, and E57
export happens in the manufacturer's processing software after scan registration \cite{cadinterop_e57}.

\subsection{Proprietary Manufacturer Formats --- FLS, PTX/PTG, TZF, ZFS}

The three biggest terrestrial laser scanner manufacturers each use their own native format.
FARO records in FLS, Leica in PTX or PTG \cite{leica_supported_formats}, and Trimble in TZF \cite{cadinterop_formats}. These
proprietary binary formats include optimizations specific to each manufacturer, ensuring maximum
preservation of the raw scan data and best performance within their own software ecosystem ---
but at the cost of depending on that vendor's ecosystem \cite{cadinterop_formats}. Zoller+Fröhlich
(Z+F) scanners use the ZFS format, which can also carry color and registration information.
Because most scanner hardware is only directly compatible with the manufacturer's own software,
users working across different scanner brands quickly run into interoperability problems
\cite{vercator2021formats}. For a project like Ballenberg --- where data has been collected over
many years, with different hardware --- understanding which proprietary formats are
present in the archive is an important first step before deciding on a conversion strategy.
For any pipeline that crosses vendor boundaries, E57 has therefore become the standard
exchange format, since the proprietary binary formats of each manufacturer cannot be read
directly by another vendor's software \cite{cadinterop_e57}.

\subsection{PLY}

PLY (Polygon File Format) was developed at Stanford University and was originally designed to
store 3D scan data from research scanners. It supports both binary and ASCII forms and can
store points, faces, colors, normals, and custom user-defined properties \cite{loc_ply}. PLY is
widely used in academic research and open-source tools like CloudCompare and Open3D, making
it a practical format for processing and algorithm development. It is also the primary input
format for 3DHOP, a web-based heritage viewer developed by the Italian research institute
ISTI-CNR \cite{potenziani20153dhop}. Unlike E57, PLY does not store scanner station data or panoramic
images, so it is better suited as a working or export format than as an archival one.

\subsection{PCD}

PCD (Point Cloud Data) is the native format of the Point Cloud Library (PCL), an open-source
framework widely used in robotics and computer vision research. PCD supports organized point
cloud datasets and uses binary memory-mapped data types, making it fast to load and save ---
an important property for real-time applications \cite{pcl_pcd}. It can store coordinates,
colors, normals, and multi-dimensional feature descriptors. For a heritage documentation
context, PCD is less relevant as a delivery or archival format, but it may appear in workflows
where automated point cloud processing algorithms from the PCL ecosystem are used, for
example in semantic segmentation tasks.

\subsection{COPC --- Cloud Optimized Point Cloud}

COPC is a newer format that extends LAZ with an internal octree-based index structure,
designed for streaming point cloud data over the web. Instead of downloading an entire file,
applications using COPC can request only the specific points needed for a given view or query,
which reduces data transfer times significantly and makes it well suited for cloud-based or
browser-based access \cite{copc_spec}. Because COPC files are valid LAZ files, any software
that can read LAZ can also open a COPC file \cite{copc_spec}, keeping backward compatibility
intact. This makes COPC a practical bridge between archival storage and interactive web
visualization --- directly relevant to the deployment goals of this project.

\subsection{EPT --- Entwine Point Tiles}

EPT is an open content organisation scheme for point cloud data designed for web-based
streaming. It is completely lossless, uses an octree-based storage structure, and stores
metadata in JSON \cite{esco2022review}. Unlike COPC, EPT stores data as a large directory of
many small individual files rather than a single file — each node in the octree becomes a
separate file on disk or in cloud storage. This design works well for streaming but creates
practical problems when moving or archiving data, since a single point cloud can result in
thousands or even millions of files. EPT is the format used natively by Potree, one of the main
visualization platforms considered in this project, making it directly relevant to the
implementation phase. In recent years COPC has emerged as a more storage-friendly alternative
that preserves the same octree streaming concept within a single file, but EPT remains widely
deployed in existing Potree-based installations.

\subsection{3D Tiles}

3D Tiles is an open standard developed by Cesium and adopted by the Open Geospatial
Consortium (OGC) for streaming large-scale heterogeneous 3D geospatial content to web
browsers and mobile applications \cite{ogc_3dtiles}. It defines a hierarchical spatial data
structure that can contain point clouds, meshes, and building models, making it more versatile
than point-cloud-only formats like EPT or COPC. Data is divided into a tree of tiles, where
each tile contains a bounded subset of the content at a given level of detail. 3D Tiles is the
native format of CesiumJS and is supported by a growing number of geospatial platforms. For
a project combining point clouds with georeferenced building models, as is the case at
Ballenberg, 3D Tiles offers a path to integrating multiple data types into a single web-based
viewer. Its relationship to the visualization platforms considered in this project is discussed
further in Section~\ref{sec:visualization}.

\subsection{HDF5}

HDF5 (Hierarchical Data Format version 5) is a general-purpose scientific data format
developed by The HDF Group, widely used in remote sensing and earth observation for storing
large, structured datasets \cite{esco2022review}. It uses a tree-like internal structure that can
hold multidimensional arrays, metadata, and multiple data types within a single file. While HDF5
is not common in terrestrial laser scanning workflows, it is worth noting in the context of the
Ballenberg project because some airborne LiDAR datasets and scientific point cloud archives use
it as their primary format. Its main limitation for interactive web use is that it was designed for
local and high-performance computing workflows and does not support efficient partial access
over a network without additional tooling.

\subsection{Format Strategy}

No single format covers every use case when working with heterogeneous 3D data from multiple
sources and capture methods. A recommended practice is to maintain a master version in a rich
format like LAS or E57, and then generate derived versions optimized for specific purposes
\cite{cadinterop_las}. For this project a layered approach makes sense: proprietary manufacturer
formats as the raw source, E57 for long-term archiving and cross-software exchange, LAZ as the
standard working format, and COPC or 3D Tiles for web deployment and interactive visualization
(discussed further in Section~\ref{sec:visualization}).

\section{Data Structuring Strategies for Heterogeneous 3D Data}
\label{sec:structuring}

Collecting 3D data over many years from different capture methods --- terrestrial laser
scanning, UAV photogrammetry, close-range photogrammetry, and mesh reconstruction --- results
in a heterogeneous archive where files differ in format, resolution, coordinate system, and
level of completeness. Before any visualization or analysis can take place, this data needs to
be organized in a way that makes it findable, understandable, and usable over the long term.
This section reviews the key strategies for structuring such an archive.

\subsection{Challenges of Heterogeneous Heritage Data}

Each capture method produces data with different characteristics. Creating a reliable 3D
dataset from heterogeneous sources presents several challenges in terms of data-gathering
pipeline, comprehensiveness, redundancy, interpretation, organization, and integration with
other heterogeneous data \cite{malik2023heterogeneous}. TLS provides high geometric accuracy
at building scale but is limited in coverage of large outdoor areas; UAV photogrammetry covers
large extents efficiently but at lower point density; close-range photogrammetry excels for
detailed object documentation. This multimodal approach --- referred to as reality-based
modeling --- advocates the strategic use of different technologies based on site
characteristics, conservation goals, and institutional resources \cite{liu2025tls}. The result,
however, is a collection of datasets that vary in spatial reference, point density, file
format, and semantic content, all of which must be reconciled before the data can be used
together.

\subsection{Spatial Indexing and Level of Detail}

The most fundamental structural challenge with large point clouds is how to organize points in
space so they can be accessed efficiently. The most widely used approach is the octree --- a
hierarchical spatial data structure that recursively subdivides 3D space into eight equal
sub-volumes. The octree method has been applied to compress, analyse, and efficiently query
point cloud data, with direct applications in cultural heritage contexts
\cite{bienkowski_use_2023}. By organizing points into an octree, systems can quickly determine
which portions of a dataset are relevant to a given view or query, loading only the necessary
data. This is the principle behind Level of Detail (LOD) rendering, where coarser
representations are shown at a distance and finer detail is loaded progressively as the user
zooms in --- the same idea behind the image pyramid familiar from GIS raster data.

Among spatial index methods, octree and KD-tree are the most frequently used for 3D point
cloud data organization. The octree performs well for uniformly distributed data, but can
become unbalanced when point density varies strongly across a scene \cite{jiang2021hybrid}.
For datasets like Ballenberg --- where dense interior scans of individual buildings sit
alongside lower-density outdoor coverage --- this non-uniform distribution is the norm. Hybrid
structures combining octrees with R-trees have been proposed to address this, offering better
nearest-neighbor search performance for uneven datasets \cite{jiang2021hybrid}. In practice,
however, the formats discussed in Section~\ref{sec:formats} (EPT, COPC, 3D Tiles) already
implement octree-based LOD internally, meaning the choice of streaming format largely
determines the spatial indexing strategy as well.

\subsection{Multi-Source Data Registration and Integration}

Before heterogeneous datasets can be stored together, they must be brought into a common
coordinate reference system through registration. The most common processing steps include
noise reduction, point cloud registration --- combining multiple point clouds from different
scan positions --- and georeferencing using GNSS or a local coordinate system to enable the
exact location of input data to be determined \cite{pepe2024integration}. For a site like
Ballenberg, where data has been captured across multiple campaigns over several years,
establishing a consistent coordinate reference across all datasets is particularly important.
Any drift or inconsistency in registration at this stage propagates through all downstream
uses.

A common integration strategy uses TLS for high-accuracy base geometry, with UAV or
photogrammetric point clouds filling in occluded or inaccessible areas that static scanners
cannot reach \cite{hassan2021scanner}. The resulting combined dataset is denser and more
complete than either source alone, but requires careful documentation of which parts of the
final cloud came from which source, at what resolution, and with what estimated accuracy. This
provenance information is as important as the geometry itself for any future conservation or
research use.

\subsection{Metadata Schemas for Heritage Documentation}

Organizing files in folders is not enough --- each dataset also needs structured metadata that
describes what was captured, how, when, and by whom. In the cultural heritage domain, the
dominant standard for this is CIDOC-CRM (Conceptual Reference Model), an ISO standard (ISO
21127) developed by the International Council of Museums. CIDOC-CRM provides definitions and a
structured framework to describe concepts, relationships, and data used in cultural heritage
documentation, and has been proposed as a common integration layer for various national and
European metadata schemas including ICCD, MIDAS, LIDO, and EU-CHIC \cite{ronchi2014metadata}.

For 3D survey data specifically, the CRMdig extension of CIDOC-CRM adds the ability to
document the provenance of digital objects --- which hardware was used, which software
processed the data, and how derived products relate to original captures. The ability to record
information about the provenance of digital objects is particularly important when the objects
are 3D digital replicas of cultural heritage assets \cite{fresa2014metadata}. Heritage Building Information Modelling (HBIM) offers a parallel approach for individual
buildings, linking 3D geometry to structured building records via the IFC data standard
\cite{murphy2009hbim, ifc_standard}, but its scope and modelling overhead make it
complementary to rather than a replacement for file-level metadata across a heterogeneous
multi-building archive. For practical metadata at the file level, a minimum set of attributes
per dataset should include: capture method, capture date, scanner or camera model, coordinate
reference system, estimated point density or resolution, file format, file size, and
processing status. This directly corresponds
to the data inventory step proposed in this project, where all existing Ballenberg datasets are
scanned and their metadata extracted as a first task.

\subsection{Archival Frameworks in Practice}

A directly relevant example is the OH3D platform, an open archival framework for sharing
cultural heritage 3D survey data spanning 36 countries and 30 UNESCO World Heritage sites
\cite{spreafico2024archival}. The platform incorporates carefully structured open data format,
authorship, and metadata frameworks reflecting the data structures of 3D capture modalities,
to establish a canonical provenance for survey data and ensure the broadest utility of these
data for posterity. Datasets are organized into projects and are searchable by various metadata
or by navigating a map. This kind of project-based, map-searchable structure --- where
each dataset is linked to a geographic location, a building or site object, and a set of
descriptive metadata --- represents a practical target for the Ballenberg archive as well.

\subsection{File Naming, Folder Structure, and Versioning}

Beyond metadata schemas, the practical organization of files on disk or in a shared storage
system matters enormously for long-term usability. A consistent folder structure --- organized
by building or site unit, then by capture campaign, then by processing stage --- reduces the
risk of data being lost or misattributed over time. File naming should encode key attributes
such as building ID, capture date, capture method, and processing level directly in the
filename, so that files remain interpretable without needing to open them or consult a separate
catalog. Version management is also important: processed derivatives should be clearly
distinguished from raw data, and any reprocessed versions should not silently overwrite earlier
outputs. These practical conventions are often underspecified in heritage documentation
workflows \cite{liu2025tls}, yet they are a prerequisite for the kind of scalable,
future-proof archive that the Ballenberg project aims to establish.


\section{Web-Based 3D Visualization Platforms}
\label{sec:visualization}

Making large heritage 3D datasets accessible to a broad audience requires more than good data
management: it requires a delivery layer that can stream massive point clouds and textured models
to a standard web browser without requiring software installation. The rapid maturation of the
WebGL standard, which exposes GPU-accelerated rendering through the browser's JavaScript
environment, has made this practical since the early 2010s. This section reviews the technical
foundation of browser-based 3D rendering, describes the main open-source platforms used in
heritage contexts, and discusses their trade-offs with reference to the requirements of a
multi-building, heterogeneous archive such as Ballenberg.

\subsection{WebGL as the Technical Foundation}

WebGL (Web Graphics Library) is a JavaScript API that provides access to the GPU directly from
within the browser, based on the OpenGL ES 2.0 specification. It became the enabling technology
for plug-in-free 3D rendering on the web and is now supported by all major browsers on desktop
and laptop platforms. WebGL 2.0 extended the API with support for more advanced shader features
and larger texture buffers. Open-source WebGL JavaScript libraries have been shown to enable the online integration of heterogeneous 3D datasets
--- including different levels of detail of buildings within a globe-based WebGIS platform ---
providing browser-accessible navigation useful for citizens, governmental bodies, and research
institutions without requiring local software~\cite{laguardia2022webgl}.

Despite its broad adoption, WebGL is reaching its limits for the scale of datasets now common in
heritage documentation. A single TLS scan campaign across a large open-air museum may produce
tens of billions of points, and WebGL's single-threaded JavaScript execution model and lack of
compute shader support constrain what can be done at rendering time. These limitations are the
main motivation behind Potree-Next (see Section~\ref{subsec:potree}) and the broader transition
toward WebGPU, the successor graphics API that provides compute shaders, multi-threading, and
a more modern GPU abstraction.

\subsection{Potree}
\label{subsec:potree}

Potree is the most widely deployed open-source WebGL point cloud viewer in the heritage and
geospatial communities. Developed at the Institute of Computer Graphics and Algorithms at TU
Wien~\cite{schuetz2016potree}, it implements a multi-resolution octree that
streams only the level of detail needed for the current view, allowing billions of points to be
explored in a standard browser. The underlying algorithm is described in detail in~\cite{schuetz2020octree}. Potree supports EPT (Entwine Point Tiles) and COPC (Cloud
Optimized Point Cloud) as input formats in addition to its own proprietary octree format, and
provides built-in tools for distance measurement, cross-section profiles, clipping volumes,
annotations, and camera animations. Its JavaScript architecture makes it straightforward to
embed in custom HTML pages and to extend with additional functionality, as demonstrated in an interactive heritage storytelling
platform for the Castello Farnese in Piacenza, combining Potree-rendered point clouds with
annotations, camera tours, and multimedia content~\cite{gaspari2023webgl}. The OH3D archival framework discussed in
Section~\ref{sec:structuring} also uses Potree as its visualization layer.

PotreeConverter, the companion preprocessing tool, underwent a complete rewrite to version 2.0.
The new converter is 10 to 50 times faster than version 1.7 on SSD storage, and critically
reduces the output from potentially millions of individual tile files to a total of three
files~\cite{potreeconverter_github}, solving a major deployment and file-system management
problem for large archives. The output format of PotreeConverter 2.0 is backward-compatible
with older Potree viewers.

\paragraph{Potree-Next: a WebGPU rewrite in progress.}
The current production version of Potree (1.8) is built on WebGL and Three.js and has
accumulated features over more than a decade of development. In October 2025, a complete rewrite of Potree under the name \emph{Potree-Next} was announced, funded
by the Austrian Netidee grant programme~\cite{potreenext_github}. Potree-Next replaces the
WebGL rendering layer with WebGPU, the next-generation browser graphics API that provides
compute shaders and a significantly more capable GPU interface. According to the project
documentation, Potree-Next is explicitly designed to take advantage of compute shaders for
point cloud processing, to eliminate legacy architectural constraints that have accumulated
over the 13-year history of the original codebase, and to target eventual deployment on Linux
and mobile devices --- two platforms where the current Potree has limited or no WebGPU support
at present. The rewrite starts with a subset of Potree 1.8 features and new capabilities that
were not previously possible; it is expected to catch up with the current feature set and
eventually replace Potree 1.8 once WebGPU achieves broad cross-platform browser
support~\cite{potreenext_github}. For projects being planned today, Potree 1.8 remains the
stable production choice, while Potree-Next should be monitored as the likely standard for
the next generation of deployments.

\subsection{CesiumJS}

CesiumJS is an open-source JavaScript library for globe-based 3D geospatial visualization,
developed by Cesium Inc. and widely used for city-scale and territorial-scale applications.
Its primary data format is 3D Tiles (OGC standard, discussed in Section~\ref{sec:formats}),
which supports point clouds, meshes, and textured buildings within a single
hierarchical tiling scheme. CesiumJS excels at multi-layer visualization: satellite imagery,
terrain models, vector features, and 3D objects can all be rendered simultaneously in a
geographically correct globe context. It supports time-dynamic simulation and 4D
visualization, making it suitable for monitoring applications where the temporal evolution
of a site is relevant.

For point clouds specifically, CesiumJS requires data to be converted to 3D Tiles format;
it does not natively read LAZ, EPT, or the Potree octree format~\cite{cesium_matom}. This
conversion step can be performed with open-source tools or through the commercial Cesium ion
cloud service. A comparison of Potree and CesiumJS for the Castello Farnese heritage site concluded that both tools are capable of hosting
rich interactive storytelling experiences with annotations, camera animations, and multimedia
integration, but that Potree offers a more direct path for point cloud deployment while
CesiumJS provides stronger georeferencing and multi-layer context~\cite{gaspari2023webgl}. CesiumJS globe-scale navigation has similarly been combined with
close-range WebGL environments to offer different levels of detail from territorial overview
to building-level inspection in a single web platform~\cite{laguardia2022webgl}.

A practical consideration for heritage institutions is that CesiumJS is open source but the
Cesium ion platform --- which handles cloud hosting, format conversion, and terrain streaming
--- operates on a subscription model. Self-hosted deployments using only open-source
components are possible but require more infrastructure setup.

\subsection{3DHOP}
\label{subsec:3dhop}

3DHOP (3D Heritage Online Presenter) is an open-source WebGL framework developed by the
Visual Computing Laboratory at ISTI-CNR (Italy), designed specifically for the cultural
heritage domain~\cite{potenziani20153dhop}. Unlike Potree and CesiumJS, which target point
clouds or geographic datasets, 3DHOP focuses on high-resolution mesh models of individual
artifacts or architectural elements. It uses the Nexus.js multiresolution format to stream
large mesh models efficiently over standard HTTP, supports PLY point clouds in single-resolution
mode without additional preprocessing, and provides ready-to-use HTML templates with hotspot
annotations, image overlay panels, and multimedia integration designed for museum-quality
presentations. The framework interconnects the 3D viewer with the surrounding HTML DOM,
making it easy to create integrated presentation schemes that combine 3D models with
photographs, text, and other multimedia content.

3DHOP's target audience ranges from museum curators with basic web experience to experienced
web developers who need a customizable point cloud and mesh viewer that is oriented to
heritage communication rather than geospatial analysis~\cite{3dhop_website}. It has been
adopted widely in the European research infrastructure for heritage science (E-RIHS), and
powers presentations at institutions including the Stavropol Museum, the Swedish Pompeii
Project, and Global Digital Heritage. At the scale of Ballenberg --- dozens of buildings
across an outdoor site --- 3DHOP is not a viable candidate platform: its object-focused
architecture and single-resolution PLY support cannot handle the volume or spatial extent
of a multi-building TLS archive. It is reviewed here as a conceptual reference, because
its approach of directly interconnecting the 3D viewer with the surrounding HTML DOM ---
enabling hotspot annotations linked to descriptive text, imagery, and conservation records
--- represents the kind of semantic linking addressed by Research Question~4. A comparable
annotation layer, adapted for site scale, is achievable within Potree or CesiumJS using
their respective extension mechanisms.

\subsection{Comparative Assessment}

For the Ballenberg context, the relevant comparison is between Potree and CesiumJS; 3DHOP
operates at object scale and is excluded from the practical selection for the reasons
described in Section~\ref{subsec:3dhop}. Potree is the stronger choice when the primary
deliverable is a navigable, measurable point cloud of a large or multi-building site; it
has the largest community of heritage deployments, supports EPT and COPC natively, and its
preprocessing and hosting pipeline is fully open-source and self-hostable. CesiumJS is the
stronger choice when geographic context, multi-layer data integration, or temporal dynamics
are important; for a museum that wants to present individual buildings within their landscape
setting and potentially link to cadastral or topographic data, CesiumJS provides the best
geospatial framework. Table~\ref{tab:webplatforms} summarizes the trade-offs between the
two candidate platforms, with Potree-Next listed as a future reference point.

\begin{table}[h]
\centering
\caption{Comparative summary of open-source web-based 3D visualization platforms for heritage.}
\label{tab:webplatforms}
\begin{tabular}{lllll}
\hline
\textbf{Platform} & \textbf{Primary data} & \textbf{Scale} & \textbf{Annotation} & \textbf{Audience} \\
\hline
Potree 1.8        & Point clouds (EPT, COPC, proprietary) & Site--region  & Basic built-in & Technical \\
CesiumJS          & 3D Tiles (mesh + points), globe layers & Region--globe & CZML / custom & Mixed \\
Potree-Next       & TBD (WebGPU, in development)           & TBD            & TBD            & TBD \\
\hline
\end{tabular}
\end{table}

\subsection{Potential Visualization Prototypes for Ballenberg}

The available Ballenberg laser scan data suggests several concrete candidates for
demonstrating the integration pipeline as mobile-optimized prototypes. The Meggen farmhouse,
documented across two distinct construction phases, could illustrate how the pipeline
represents temporal layers within a single building --- showing visitors the evolution of
the structure without requiring navigation of the full point cloud. The settlement history
dataset covering multiple relocated structures could underpin a site-wide narrative view
contextualizing individual buildings within the broader collection. The Villa Braham, with
its documented structural system, lends itself to annotated cross-sections illustrating
both vertical and horizontal load distribution --- a use case relevant to conservation
staff and educational audiences alike. In each case the intended interaction model is a
focused, single-aspect presentation rather than an open-ended explorer, which directly
shapes the annotation and interface requirements for whichever platform is ultimately
selected. Mobile optimization is a baseline requirement for all three, given that
on-site visitors are the primary intended audience.

\subsection{Limitations and Open Challenges}

Several open problems limit current web-based heritage visualization. First, mobile device
support is inconsistent: WebGL performance on smartphones and tablets varies widely by
manufacturer and operating system, and the current Potree 1.8 does not run reliably on
mobile browsers. Potree-Next explicitly targets mobile support as a goal, but this will
only be achievable once WebGPU is broadly supported on mobile platforms. Second, semantic
metadata integration remains shallow in both candidate platforms: annotations can
display text, images, and links, but there is no standardized mechanism for displaying
CIDOC-CRM property sets as structured, queryable metadata within the point cloud
viewer. The connection between the visual geometry and the semantic documentation layer
lives outside the viewer, typically in an external database or documentation system. Third, combining
point clouds with meshes, orthophotos, and 3D models in a single coherent scene requires
manual data preparation and format conversion work that is not yet automated in any of the
reviewed platforms. For an archive of the scale and heterogeneity of Ballenberg, building
a web visualization layer therefore involves not only selecting a platform but designing
an integration pipeline that aligns all data types in a common coordinate system,
converts each to the appropriate streaming format, and links the visual output to the
metadata layer.
\section{Synthesis and Research Gaps}
\label{sec:synthesis}

The four themes reviewed in this report --- point cloud formats and standards, data structuring
strategies, semantic metadata for heritage documentation, and web-based visualization platforms ---
are not independent topics. They form a pipeline in which choices made at one stage propagate
as constraints or opportunities at every subsequent stage. This section synthesizes the
cross-cutting dependencies between themes, identifies the research gaps that emerge from the
review, and positions the Ballenberg project within that landscape.

\subsection{Cross-Cutting Dependencies}

The choice of point cloud format directly determines what is possible at the visualization and
metadata stages. Proprietary manufacturer formats (FLS, PTX, TZF) confine data to a single
software ecosystem and cannot be served to a web viewer without conversion. Archival formats
such as E57 preserve scanner station data and panoramic imagery but require intermediate
conversion before any web delivery is possible. The streaming formats COPC and 3D Tiles are
the only ones that can be loaded directly into a browser-based viewer without preprocessing
on the server side, and each is tightly coupled to a specific viewer ecosystem: COPC to Potree
and EPT-compatible viewers, 3D Tiles to CesiumJS. This means that a format strategy and a
visualization platform strategy must be decided together, not sequentially.

Data structuring and spatial indexing in turn determine how well semantic metadata can be
attached and maintained over time. An archive organized into octree-indexed tiles (EPT, COPC,
or 3D Tiles) exposes a spatially coherent hierarchy that can be used as an anchor for
point-level or region-level attribute annotation. An unstructured flat collection of LAZ files
provides no such anchor: metadata must live in a separate database with only file-level
granularity. The metadata schemas reviewed in Section~\ref{sec:structuring} --- CIDOC-CRM
and its CRMdig extension --- operate at the dataset and object level, not at the point level,
so even well-structured archival metadata does not fully bridge the gap between a record in a
heritage database and a specific cluster of points in a viewer. Closing this gap at the visualization level requires region-level annotations keyed to
external documentation records, or a spatially indexed metadata layer that can be queried
directly from within the viewer --- neither of which is provided out of the box by any of
the platforms reviewed.


Web visualization platforms inherit all of these constraints. Potree is currently the most
mature platform for point cloud delivery at site scale, but its annotation layer is not
semantically structured: hotspots contain free text and media, not ontology-aligned property
sets. CesiumJS provides richer geospatial context and supports 3D Tiles as a container for
both geometry and attributes, but the pipeline from raw TLS data to a structured CesiumJS
scene requires more preprocessing steps than a Potree deployment. The upcoming
Potree-Next rewrite may shift this landscape once WebGPU support matures and compute shaders
become available for in-browser semantic filtering, but this remains a development-stage
prospect rather than a deployable solution at the time of writing.

\subsection{Research Gaps}

\paragraph{Gap 1: No standardized heterogeneous-to-web pipeline for heritage archives.}
The literature documents many individual components of a survey-to-web pipeline --- format
converters, octree generators, viewer libraries --- but no study has established a validated,
end-to-end workflow that takes heterogeneous input (TLS + UAV photogrammetry + close-range
photogrammetry acquired across multiple campaigns) and produces a unified, web-accessible
output with consistent coordinate systems, consistent metadata, and manageable hosting
requirements. Existing studies address multi-source data integration for individual buildings~\cite{pepe2024integration, hassan2021scanner},
but the challenge of scaling this to a multi-building, multi-year archive --- where new
campaigns must be ingested without disrupting existing data --- has not been systematically
addressed. For Ballenberg, with its dozens of buildings documented over years by different
teams and instruments, this is the central practical challenge.

\paragraph{Gap 2: Lightweight metadata standards for point-cloud-first archives.}
In large heritage archives, point cloud datasets are often stored with minimal structured
annotation beyond basic file-level descriptors such as capture date and instrument. CIDOC-CRM
provides a rich ontological vocabulary but operates at an abstraction level above the file,
and its implementation requires specialist knowledge. A practical gap exists for a lightweight,
point-cloud-compatible metadata schema that can be applied consistently at dataset level across
a heterogeneous archive and that is interoperable with CIDOC-CRM for future enrichment.
The minimum metadata field set proposed in Section~\ref{sec:structuring} is a starting point,
but has not been validated against the specific documentation and conservation workflows of
an open-air museum.

\paragraph{Gap 3: Platform selection criteria for multi-building open-air museum contexts.}
The most rigorous published head-to-head evaluation of Potree and CesiumJS~\cite{gaspari2023webgl} focuses on a single monumental site and
evaluates the platforms primarily from the perspective of storytelling and visual experience.
No study has systematically compared web visualization platforms for the specific combination
of requirements that characterizes an open-air museum: dozens of buildings at varying scales,
heterogeneous data types per building, multiple simultaneous user personas (conservation
specialist, museum curator, general public), and the need to link the visual output to
structured documentation databases. The trade-offs between Potree's point-cloud-native
streaming and CesiumJS's geospatial multi-layer capabilities have not been evaluated against
this requirement profile.

\paragraph{Gap 4: Semantic linking between the web viewer and the documentation layer.}
Across all platforms reviewed, semantic metadata lives outside the viewer and is connected
to the 3D geometry only through manual annotation. There is no standardized API or data
exchange format that allows a CIDOC-CRM property set to be queried from within a
point cloud viewer, or that would allow a user to click on a building element in Potree
and retrieve structured conservation records from a linked database. One promising direction
is the integration of Linked Open Data infrastructure --- where each documented building
element or survey record is a URI that can be queried from within a viewer --- but this
has not been demonstrated at operational scale for a real heritage archive with a streaming
point cloud viewer.

\subsection{Positioning the Ballenberg Project}

The Ballenberg Open-Air Museum presents a use case that sits at the intersection of all
four identified gaps. Its archive is heterogeneous by construction: buildings from different
Swiss regions and construction periods have been documented by different survey teams over
multiple years, using different instruments and with varying levels of completeness. No single
format, metadata schema, or visualization platform has yet been applied consistently across
the full archive. The project therefore has the opportunity to contribute to the state of
the art not only by building an operational pipeline but by documenting the decisions,
trade-offs, and failure modes encountered when applying current best practices to a real,
large-scale, resource-constrained institutional archive.

The research questions stated at the outset of this report map directly onto the four gaps:
the pipeline question (Gap~1) addresses how heterogeneous survey products can be standardized
for web deployment; the metadata question (Gap~2) addresses what schema is sufficient for
scalable, interoperable management of a point-cloud-first archive; the platform question (Gap~3) addresses
which visualization tool best serves the museum's mixed user base; and the semantic linking
question (Gap~4) addresses how structured documentation can be made accessible from within
the visual interface. The synthesis presented in this section suggests that a pragmatic,
layered approach --- archival quality first in E57 and LAZ, structured metadata at dataset
level aligned with CIDOC-CRM, COPC or 3D Tiles for streaming, and an annotation layer in
Potree or CesiumJS linked to an external documentation database --- offers the most viable
path toward a future-proof, web-accessible, and institutionally maintainable archive for
the Ballenberg collection.

\section*{Project Timeline}
\label{sec:timeline}

The project is structured into five workpackages distributed across the spring semester 2026,
which began in calendar week 7. The first workpackage, the State of the Art review, was
conducted in weeks 8 to 11 and forms the foundation for all subsequent decisions regarding
formats, data structuring, semantic metadata, and visualization platform selection. In parallel,
the report writing workpackage spans weeks 8 to 11 for the STAR phase and resumes in weeks 19
to 23 for the final project report. The second workpackage covers the acquisition and preparation
of ground truth data from the Ballenberg archive, including format auditing, coordinate system
harmonization, and metadata inventory of existing datasets, and is planned for weeks 12 to 14.
The third and central workpackage involves the design and implementation of the integration
pipeline: format conversion to archival and streaming targets, spatial indexing, metadata
schema definition, and deployment of a web-based visualization environment; this phase runs
from week 13 to week 17 and overlaps deliberately with the tail of the data acquisition phase
to allow iterative refinement as real data is ingested. The fourth workpackage covers the
evaluation of results, including performance testing of the web viewer, assessment of metadata
completeness, and comparison of platform trade-offs against the defined research questions;
this is scheduled for weeks 17 to 20. The timetable is summarized in Table~\ref{tab:timeline}.

% Required in preamble: \usepackage{tikz}
\begin{figure}[h]
\centering
\caption{Project timetable by calendar week (semester start: CW\,7).}
\label{tab:timeline}
\begin{tikzpicture}[x=0.55cm, y=0.6cm]

  % --- settings ---
  \def\rows{5}
  \def\cols{16}
  \definecolor{ganttblue}{RGB}{26,35,126}

  % column labels (CW 8 to 23)
  \foreach \i in {1,...,16} {
    \pgfmathtruncatemacro{\cw}{\i+7}
    \node[font=\small\bfseries, anchor=south] at (\i-0.5, \rows+0.15) {\cw};
  }
  % CW header
  \node[font=\small\bfseries] at (7.5, \rows+0.85) {CW};

  % row labels
  \node[font=\small, anchor=east] at (-0.1, 4.5) {STAR};
  \node[font=\small, anchor=east] at (-0.1, 3.5) {Acquire GT Data};
  \node[font=\small, anchor=east] at (-0.1, 2.5) {Implement Pipeline};
  \node[font=\small, anchor=east] at (-0.1, 1.5) {Evaluate Results};
  \node[font=\small, anchor=east] at (-0.1, 0.5) {Write Report};

  % grid
  \draw[gray!40, thin] (0,0) grid[xstep=1, ystep=1] (16,5);
  \draw[black, thin] (0,0) rectangle (16,5);
  \draw[black, thin] (0,5) -- (16,5);

  % bars: {start col (0-indexed), row bottom, width, height}
  % STAR: CW8-11 = cols 0-3
  \fill[ganttblue] (0, 4.15) rectangle (4, 4.85);
  % Acquire GT Data: CW12-14 = cols 4-6
  \fill[ganttblue] (4, 3.15) rectangle (7, 3.85);
  % Implement Pipeline: CW13-17 = cols 5-9
  \fill[ganttblue] (5, 2.15) rectangle (10, 2.85);
  % Evaluate Results: CW17-20 = cols 9-12
  \fill[ganttblue] (9, 1.15) rectangle (13, 1.85);
  % Write Report: CW8-11 = cols 0-3, CW19-23 = cols 11-15
  \fill[ganttblue] (0,  0.15) rectangle (4,  0.85);
  \fill[ganttblue] (11, 0.15) rectangle (16, 0.85);

\end{tikzpicture}
\end{figure}



    \input{chapters_draft}

    \newpage
    \pagenumbering{Roman}
    \setcounter{page}{4}

    %%%%%%%%%%%%%%%%%%%%%%%%%%%%%%%%%%%%%%%%%%%%%%%%%%%%%%%%%%%%%%%%%%%%%%%%%%%%
	% Bibliografie
	%%%%%%%%%%%%%%%%%%%%%%%%%%%%%%%%%%%%%%%%%%%%%%%%%%%%%%%%%%%%%%%%%%%%%%%%%%%%
    
    \printbibliography
    \setcounter{page}{5}
    \newpage


    %%%%%%%%%%%%%%%%%%%%%%%%%%%%%%%%%%%%%%%%%%%%%%%%%%%%%%%%%%%%%%%%%%%%%%%%%%%%
	% Anhang
	%%%%%%%%%%%%%%%%%%%%%%%%%%%%%%%%%%%%%%%%%%%%%%%%%%%%%%%%%%%%%%%%%%%%%%%%%%%%
    \listoffigures
    \listoftables
    \setcounter{page}{6}
    
    \appendix{}


	
\end{document}
